% Journal of Mathematical Sciences — National Mathematical Centre, Abuja
% Vol. 8 Issue 1. 12 pt, A4, English. First page carries title, authors, affiliations,
% emails, abstract (<= 250 words, single-spaced), subject classification, keywords.
% Body order is fixed: Introduction, Methods, Results, Discussion, Conclusions,
% Declaration, References (APA 7th). Do not reorder or rename these sections.
% Submission: Editor, Professor James I. Ajie <nmmcmathsprog@gmail.com>. Deadline 2026-12-30.

\documentclass[12pt,a4paper]{article}

\usepackage[T1]{fontenc}
\usepackage[utf8]{inputenc}
\usepackage[margin=2.5cm]{geometry}
\usepackage{amsmath,amssymb,amsthm}
\usepackage{graphicx}
\usepackage{booktabs}
\usepackage{setspace}
\usepackage{hyperref}

% Long unbreakable strings (file paths in \texttt) cannot be hyphenated, so a
% paragraph carrying one runs into the margin. This lets TeX stretch spaces
% instead of overrunning the text block.
\setlength{\emergencystretch}{3em}

\newtheorem{proposition}{Proposition}
\newtheorem{lemma}[proposition]{Lemma}
\newtheorem{corollary}[proposition]{Corollary}
\theoremstyle{definition}
\newtheorem{definition}[proposition]{Definition}
\newtheorem{assumption}[proposition]{Assumption}
\theoremstyle{remark}
\newtheorem{remark}[proposition]{Remark}

% Missing-number marker. Every \todo must be empty before submission.
\newcommand{\todo}[1]{\textbf{[TODO: #1]}}

% The abstract must be 12 pt: the article class default sets it in \small, which
% would be a formatting deviation from the poster's contract.
\renewenvironment{abstract}{%
  \par\vspace{0.8em}%
  \begin{center}\textbf{Abstract}\end{center}%
  \normalsize\noindent\ignorespaces
}{%
  \par\vspace{0.8em}%
}

\title{Sharp bounds for the number of areas above a threshold, with a
spatial cross-validation protocol: severe household food insecurity in
Nigerian local government areas}

% The journal requires the corresponding author to be listed FIRST in a multi-author paper.
\author{%
  Jerry Bannister Zachary\\
  Abubakar Tafawa Balewa University\\
  \texttt{jzbannister@atbu.edu.ng}
}
\date{}

\begin{document}
\maketitle

\begin{abstract}
Food-insecurity decisions are made for local government areas (LGAs), whereas estimates are commonly reported for states. The number of areas above a policy threshold is a separate posterior quantity: its variance depends on pairwise exceedance probabilities, which marginal probabilities alone cannot provide. Local allocation therefore requires joint uncertainty about the count, not only a separate risk estimate for each area.

We analyse the 2024 Nigeria Demographic and Health Survey across 774 LGAs, where severe household food insecurity was 29.1\% nationally. A Bayesian hierarchical spatial model includes an unstructured area effect and an independently fitted cluster effect. At a 50\% prevalence threshold, the fitted expected count was 115 LGAs, with a 95\% interval of 101--130. An independence calculation would understate the count's standard deviation by 14\%; the marginals alone permit a standard deviation from 0.156 to 74.11, so the interval reflects the model's dependence structure rather than the marginals alone. The design leaves 120 LGAs without a sampled cluster, covering 11.04\% of the population. Consequently, the identified count spans 87--231, and 120 of the 144-area width rests on assumption.

With pass rules fixed before fitting, spatial cross-validation coverage was 0.890 at the exposure faced by unclustered LGAs and 0.775 when an entire region was withheld. These tests assess area-level prediction, not the count interval, which they do not validate. What transfers beyond this application are the sharp bounds and the pre-declared protocol, not a new estimator.

\end{abstract}

\vspace{0.5em}
\noindent\textbf{Subject classification (MSC 2020):} 62F15 (primary); 60E15; 62H11; 62D05; 62P25.

\vspace{0.3em}
\noindent\textbf{Keywords:} small-area estimation; partial identification; Fr\'echet--Hoeffding bounds;
comonotonicity; spatial cross-validation; food insecurity; Nigeria.

% ---------------------------------------------------------------------------
% Fixed section order. Each section lives in its own file and is included when
% it exists, so the document compiles while sections are still being drafted.
% ---------------------------------------------------------------------------
\IfFileExists{introduction.tex}{\section{Introduction}

Severe household food insecurity is commonly summarized for Nigeria at state level, but a state average cannot show where risk is concentrated within a state. The weighted national estimate from the 2024 Nigeria Demographic and Health Survey is 29.10\%. Decisions directed to local government areas (LGAs) require estimates below state level, together with a clear account of how much those estimates and any resulting area count remain uncertain. The relevant target is not only the prevalence in each LGA, but also the number of LGAs whose prevalence crosses a stated threshold.

Household severity and area status are separate decisions. The household outcome defines whether a respondent meets the survey's severe-food-insecurity definition; the LGA threshold classifies an area according to its prevalence. A household outcome cannot by itself identify whether an LGA crosses that line, and an area estimate is not a direct count of affected households. Keeping these levels distinct prevents a household measure, an area prevalence, and a count of places from being treated as the same quantity.

Spatial models of household food-insecurity items provide one relevant precedent. Chacón-Montalván et al. (2025; \href{https://doi.org/10.1214/25-aoas2072}{DOI}) fit a spatial item-factor model to data from Ipixuna in the Brazilian Amazon and map three latent factor scores derived from 18 modified food-insecurity items as a continuous surface. They also report pointwise exceedance probabilities, \(\Pr\{\theta_k(s)>0\}\), use a 0.7 probability rule to label hotspots, and count the resulting hotspot regions: four in the full data, one in the dry season, and four in the rainy season. These results make clear that exceedance probabilities and hotspot counts are already used in spatial food-insecurity work. The mapped quantity, however, is a latent factor score at a point, not an LGA prevalence. Ipixuna is one municipality, and the finest administrative unit discussed is eight census sectors; the paper does not report estimates for administrative areas.

A related article by Moura, Pereira, and Silva (2025), \emph{A Beta--Beta Prime Model for Rates and their Precision for Small Area Estimation: An Application to Brazilian Food Insecurity Index}, is closed access and could not be read in full; we make no claim about its contents (\href{https://doi.org/10.1093/jssam/smaf002}{DOI}).

Other spatial methods address different outcomes and prediction tasks. Michal et al. (2024; \href{https://doi.org/10.1177/09622802241293776}{DOI}) extend the Besag--York--Mollié (BYM2) model with a scale mixture that allows latent variance to vary across areas, propose two priors for the mixture scale, and report 23 of 160 Rio de Janeiro districts as potential outliers in an analysis of Zika cases. MacBride et al. (2025; \href{https://doi.org/10.1214/24-aoas1969}{DOI}) describe SPAR-Forest, which iteratively refits a random forest with a Bayesian conditional autoregressive (CAR) or spatial autoregressive (SAR) model. Their abstract reports better performance than CAR/SAR models, random forests, and geographical random forests on Scottish property prices and in simulations. These studies provide precedents for heavy-tailed area effects and spatial prediction, but they do not change the estimand considered here: the prevalence of severe food insecurity among households in each Nigerian LGA.

Area-level estimation for food insecurity and nutrition is an active practice in settings where surveys are designed for larger units. Chandra and Verma (2022) estimate food-insecurity prevalence for districts of Uttar Pradesh from a survey intended to be reliable at state level; Yilema et al. (2022) combine survey and census data to estimate undernutrition among children under five in Ethiopian districts; Wulandari, Darmawan and Zamzanah (2025) estimate stunting prevalence at subdistrict level in Indonesia with a hierarchical Bayesian model and conditional autoregressive spatial priors; and Santillana et al. (2023) estimate a survey proportion for every municipality in Mexico with the Fay--Herriot area-level model. These papers establish that borrowing strength across areas for nutrition and food insecurity is established practice. None of them reports an interval for the number of areas above a threshold.

The distinction between a spatial surface and an area prevalence matters for both comparison and decisions. A probability that a latent factor is positive is not the probability that more than a chosen share of households in a local government area is severely food insecure. Likewise, a list of areas with high pointwise probabilities does not give the uncertainty in the number of areas above a prevalence threshold. The available evidence must therefore be assessed against the actual decision quantity rather than against a superficially similar map or hotspot label.

Within Nigeria, published estimates exist at two levels. Cross-sectional surveys report household food insecurity for single local government areas, for example Egor in Edo State (Omuemu, Otasowie and Onyiriuka, 2012) and Akinyele in Oyo State (Kayode and Alabi, 2025). At a wider scale, Jerumeh (2024) maps household food insecurity across the 109 senatorial districts from the 2018/19 Nigeria Living Standard Survey, classifies 28 districts as hotspots and 19 as coldspots, and reports spatial dependence between neighbouring districts in a spatial regression; only its published abstract was available to us, its full text not being open access. The senatorial district is an administrative unit above the local government area. In this work we found no interval for an area estimate and no uncertainty statement for the number of areas above a line. Experience-based food insecurity in Nigeria has also been analysed against agricultural productivity (Villacis, Mayorga and Mishra, 2022), and a reproducible machine-learning pipeline predicts household-level food insecurity from the 2024 survey used here, with validation on cluster-grouped splits (Osung, 2026); we found no area-level prevalence estimates for the local government areas in either. Our claim is therefore narrower than the absence of Nigerian estimates: we found no estimates for the national set of 774 local government areas, and no published interval for the count of areas above a threshold.

This scope does not imply that exceedance probabilities are absent from prior work. Chacón-Montalván et al. report them and use a probability rule to count hotspots. Their threshold is zero on a latent factor scale, not a threshold on prevalence, and their hotspot counts are descriptive, without an interval for the count. Their paper reports model-fit comparisons and Markov chain Monte Carlo (MCMC) diagnostics, but no out-of-sample validation. The gaps we can verify in the closest paper concern the geographic unit, a prevalence threshold, uncertainty for the count, and out-of-sample performance.

The uncertainty in a count is not a technical detail. Let \(\theta_s\) denote the severe-food-insecurity prevalence in LGA \(s\), let \(c\) be a chosen threshold, and define \(I_s(c)=\mathbf{1}\{\theta_s>c\}\). The count of areas above the threshold is \(C(c)=\sum_{s=1}^{S}I_s(c)\), where \(S\) is the number of LGAs. If \(p_s=\Pr\{I_s(c)=1\}\) and \(P_{st}=\Pr\{I_s(c)=1,I_t(c)=1\}\), then
\[
  \operatorname{E}\{C(c)\}=\sum_{s=1}^{S}p_s,\qquad
  \operatorname{Var}\{C(c)\}
  =\sum_{s=1}^{S}p_s(1-p_s)
    +2\sum_{s<t}\left(P_{st}-p_sp_t\right).
\]
The mean is fixed by the area-specific probabilities because expectation is additive. The variance is not: it depends on whether exceedances occur together or in different places, information that the marginal probabilities do not contain. In this application, the standard deviation of the count can range from 0.16 to 74.11 under joint distributions that all have the same marginals. This range is not a confidence interval around one fitted estimate. It describes how much the uncertainty in the count can vary when only the marginal probabilities are known. An interval based on one fitted spatial dependence structure is a different, model-dependent quantity, and area-level prediction checks do not by themselves validate that joint count interval. A number with no interval, or with an interval that has never been checked, is not usable for targeting.

This distinction also affects how spatial validation should be interpreted. Prediction for one unsampled LGA while nearby LGAs remain in the fit is not the same task as prediction for an entire region whose local observations and neighbours are withheld. A validation result should match the intended use, and a pass for individual area estimates should not be presented as evidence that the interval for the joint count is calibrated. The paper therefore treats the marginal bounds, the model-based count distribution, and the spatial validation results as separate pieces of evidence.

Cross-validation is a design question, not a single score. Withheld LGAs may be isolated or grouped, and neighbouring observations may remain available or disappear with them. Those designs represent different prediction settings and answer different questions about how far the estimates can be carried. We therefore define the fold structure and pass rules before comparing predictions, rather than treating one split as a universal test of spatial performance. The aim is to identify which prediction tasks the fitted estimates support and which they do not.

We use established probability results and spatial models to examine these questions in the 2024 Nigerian application. The paper makes four contributions:
\begin{enumerate}
  \item We derive sharp bounds on the variance of the number of LGAs above a prevalence threshold when only the marginal exceedance probabilities are specified.
  \item We set out a spatial cross-validation protocol with pass rules declared in advance and apply it in four designs, distinguishing prediction for one area among sampled neighbours from prediction for a whole region without local observations.
  \item We give a partial-identification statement for LGAs with no sampled cluster, separating what the survey identifies from what must remain unknown.
  \item We apply the analysis to severe household food insecurity in Nigeria and find that the environmental covariates tested carry no measurable signal at LGA level.
\end{enumerate}
The contribution is an application and evaluation of existing ideas, not a new estimator. The count bounds clarify what marginal probabilities alone can establish; the holdout designs assess area prediction under distinct spatial information settings; and the partial-identification statement prevents model predictions for unclustered areas from being mistaken for direct observations. Together, these components keep a fitted map, a count interval, and the survey's limits visible as different claims.

The remainder of the paper is organized as follows. Methods defines the outcome, household and area model, count functional, partial-identification bounds, and spatial validation protocol, and presents the supporting propositions. Results reports the descriptive estimates, area summaries, count uncertainty, bounds for areas without clusters, validation outcomes, and covariate analysis. Discussion considers the interpretation and limits of these results. Conclusions summarize the findings, followed by the Declaration and References.
}{}
\IfFileExists{methods.tex}{% Methods section.  The TODO command keeps unresolved specifications visible
% while allowing this fragment to compile under a bare article wrapper.
\providecommand{\todo}[1]{\textbf{[TODO: #1]}}

\section{Methods}

\subsection{Data and setting}

We used the Nigeria Demographic and Health Survey 2024 to estimate severe household food insecurity at local government area (LGA) level. The analysis unit was the household, while the target areas were the country's 774 LGAs. The analysis retained all 39,877 households in the extract, with no missing covariates in the analysed set. These households came from 1,380 sampled clusters in 37 jurisdictions, comprising the 36 states and the Federal Capital Territory. The survey is clustered and stratified, and the household records carry sampling weights. Those design features matter because the fitted likelihood is unweighted; we therefore distinguish the sample quantity estimated by the model from the design-weighted national prevalence.

State summaries are not the target of the spatial model. The direct severe-food-insecurity estimates range from 4.6\% in Niger to 85.0\% in Yobe, with a median of 23.5\%. A state average cannot show how risk is distributed among the LGAs within that state. The design work for this study found that direct state estimates were already sufficiently precise for state-level summaries, leaving the within-state distribution as the unresolved problem. A Fay--Herriot calculation on the 37 states gives a median shrinkage weight of 0.987, so the model moves the direct estimate by 0.0026 on average and by at most 0.020: at state level there is almost nothing left to borrow. The LGA analysis is intended to describe this finer geography, not to replace the state estimates with a more elaborate estimate of the same state average.

The cluster coordinates were assigned to LGA polygons by spatial overlay. We checked each polygon's state against the survey's jurisdiction label. Where they disagreed, the survey jurisdiction was treated as authoritative; 15 clusters, or 1.1\% of the total, had this mismatch. The LGA graph was built using queen contiguity: two areas are neighbours when their polygons share a boundary point, whether along an edge or at a vertex. The resulting graph has 774 nodes, 2,230 edges, and one connected component. This graph defines the structured part of the area effect and also supplies the neighbourhoods used to design the calibration tests.

The unweighted prevalence in the analysis extract was 27.74\%, compared with a design-weighted national prevalence of 29.10\%. We did not use weights in the likelihood and do not interpret the fitted national mean as the design-weighted population estimate. The difference is reported here because the spatial model estimates the unweighted household probability. The random cluster effect described below represents residual dependence among sampled households in a cluster; it does not correct selection probabilities or reproduce the survey's stratification and weighting.

\subsection{Outcome}

For household \(h\), let \(r_h\) denote the survey office's Rasch probability for severe food insecurity. We defined the binary response as
\[
  Y_h = \mathbf{1}\{r_h \geq 0.50\}.
\]
Thus, the modelled response is the survey office's severity construct thresholded at 0.50. We did not define severe food insecurity by summing affirmed items and applying a raw-count cut-off. A raw count gives each item the same contribution, whereas the Rasch probability represents the survey office's calibrated severity measure. The choice of outcome is therefore a choice of construct, not just a convenient recoding of the same measure.

We kept the household outcome definition fixed and examined threshold sensitivity at the area level. For a crisis threshold \(c=0.50\), the area count is the number of LGAs whose latent prevalence exceeds \(c\). Repeating the calculation at \(c=0.40\) and \(c=0.60\) gives posterior expected counts of 196.28 and 61.28, respectively. These alternatives change the area-level crisis line; they do not reclassify households or substitute a raw item count for the survey office's outcome.

\subsection{Household and area model}

We fitted a household-level binomial model. Conditional on its probability \(\pi_h\), each response follows
\[
  Y_h \mid \pi_h \sim \operatorname{Bernoulli}(\pi_h),
  \qquad
  \operatorname{logit}(\pi_h)
  = \beta_0 + \mathbf{x}_h^\mathsf{T}\boldsymbol{\beta}
  + \mathbf{w}_{s(h)}^\mathsf{T}\boldsymbol{\gamma}
  + u_{s(h)} + b_{g(h)} ,
\]
where \(s(h)\) is the LGA containing household \(h\), \(g(h)\) is its sampled cluster, \(\mathbf{x}_h\) is its household-covariate vector, and \(\mathbf{w}_{s(h)}\) is the vector of area covariates. The household covariates were wealth (a standardised asset index), household size, an urban indicator, whether the head is male, the head's age centred and its square, two education indicators for the head (secondary; higher), an improved-water indicator, and an improved-sanitation indicator: ten design columns for eight covariates. Improved sanitation is coded zero, not missing, for a household with no toilet facility; treating those blanks as missing would have deleted the poorest quarter of the sample.

The area effect \(u_s\) follows the BYM2 construction on the queen-contiguity graph. In this parameterisation,
\[
  u_s = \sigma_u\!\left(
    \sqrt{\phi}\,u_s^{(S)}
    + \sqrt{1-\phi}\,u_s^{(I)}
  \right),
\]
where \(u^{(S)}\) is the graph-structured component, scaled to the graph, and \(u^{(I)}_s\) is an independent standard normal component. The mixing parameter \(\phi\) controls the relative contribution of the structured component, and \(\sigma_u\) sets the overall scale. We centre the structured component to distinguish it from the intercept. The cluster effect is an independent Gaussian random intercept,
\[
  b_g \stackrel{\mathrm{iid}}{\sim} N(0,\sigma_b^2),
\]
shared by households sampled in cluster \(g\). This effect represents residual within-cluster similarity after conditioning on the covariates and LGA effect.

Eight candidate area covariates were interpolated by inverse-distance weighting, using the three nearest of the 1,380 sampled clusters; population density was computed from the 2020 population projections and polygon area rather than interpolated. Their cross-validated interpolation correlations were 0.999 for aridity, 0.999 for precipitation, 0.977 for elevation, 0.923 for malaria, 0.920 for vegetation, 0.614 for travel time, 0.535 for livestock, and 0.522 for nightlights. Selection was mechanical and fixed before any fit: drop a candidate whose cross-validated interpolation correlation is below 0.85; then drop any candidate correlated above 0.8 with one already kept, the earlier in a stated order winning; then keep at most six. Five survived: population density, aridity, vegetation, elevation and malaria.

\textbf{No area covariate enters the primary model.} Every result reported in this paper comes from a fit with household covariates only. The five area covariates were added in a separate specification and estimated effects indistinguishable from zero, in nineteen of twenty coefficient intervals; the primary results therefore exclude them, and that exercise is reported as a negative finding rather than as a specification. Three of the eight candidates could not be recovered at this spatial resolution at all, and that is also reported.

The area-specific prevalence \(\theta_s\) is the model-implied probability at the LGA covariate profile, averaged over the cluster-effect distribution rather than evaluated at a zero cluster effect. In posterior draw \(m\), we calculate it as
\[
  \theta_s^{(m)}
  = \int \operatorname{logit}^{-1}\!\left(
       \beta_0^{(m)}
       + \bar{\mathbf{x}}_s^\mathsf{T}\boldsymbol{\beta}^{(m)}
       + \mathbf{w}_s^\mathsf{T}\boldsymbol{\gamma}^{(m)}
       + u_s^{(m)} + b
    \right)\,dF_b^{(m)}(b),
\]
where \(\bar{\mathbf{x}}_s\) is the household-covariate profile used for the LGA prediction and \(F_b^{(m)}\) is the draw-specific Gaussian distribution for the cluster effect. We integrated over that distribution numerically. The profile \(\bar{\mathbf{x}}_s\) is the LGA's mean household covariate profile, computed from the full household file, and the cluster effect was marginalised by Gauss--Hermite quadrature.

The likelihood deliberately omits the survey weights. Consequently, its estimand is the unweighted probability in the analysed sample, not a design-weighted population prevalence. The LGA random effect and cluster random effect do not make the likelihood design-weighted. Nor were the held-out covariates re-estimated using only the training folds. They were supplied as known at prediction time. This choice gives the model information that may not be available for a genuinely unsampled LGA, especially for household covariates, and the validation results should be read as conditional on that information being available. They do not measure uncertainty from estimating or transporting those covariates.

The BYM2 precision used a penalised-complexity prior on the log precision, set so that the prior probability of a standard deviation above 1 is 0.01; the cluster precision used INLA's default penalised-complexity prior for an IID effect; the fixed effects used INLA's default Gaussian prior. We do not infer these settings from posterior summaries.

\subsection{Count of areas above a threshold}

For an area-level crisis threshold \(c\), let \(I_s=\mathbf{1}\{\theta_s>c\}\) and define
\[
  C(c)=\sum_{s=1}^{774} I_s .
\]
The target is the posterior distribution of this count, not only a list of separate LGA estimates. For each joint posterior draw, we calculate every \(\theta_s\), apply the threshold, and sum the resulting indicators. We use those same draws to obtain the posterior mean, interval, and variance of \(C(c)\); these functionals require no additional likelihood evaluations.

Write \(p_s=\Pr(\theta_s>c)\) and \(P_{st}=\Pr(\theta_s>c,\theta_t>c)\). Proposition 1 in the theory section gives
\[
  \operatorname{E}[C(c)] = \sum_s p_s,\qquad
  \operatorname{Var}\{C(c)\}
  = \sum_s p_s(1-p_s)
    + 2\sum_{s<t}\left(P_{st}-p_sp_t\right).
\]
The mean is marginal because expectation is additive, so the area-specific exceedance probabilities determine it. The variance also depends on pairwise joint exceedance probabilities, so it is not determined by the marginals alone. We report the posterior interval as a model-based summary of the joint draws. Proposition 1 and the accompanying theory section provide the formal results; we do not repeat their proofs here.

\subsection{Partial identification for areas without clusters}

The sample does not contain a cluster in every LGA. We let \(\mathcal{O}\) be the set of LGAs with at least one sampled cluster and \(\mathcal{U}\) the set with none. The observed set contains 654 LGAs, while 120 LGAs are unobserved. The population share in \(\mathcal{O}\) is \(p=0.88963\), leaving 11.04\% of the population in \(\mathcal{U}\). These population shares, rather than the fraction of geographic units, define the bounds for the share of people living in a crisis area.

For each posterior draw, the observed-area count is
\[
  C_{\mathcal{O}}^{(m)}(c)
  = \sum_{s\in\mathcal{O}}
      \mathbf{1}\{\theta_s^{(m)}>c\}.
\]
We place no restriction on the prevalence in an LGA without a cluster. Its contribution to the full count may therefore be zero or one, giving the draw-specific identified set
\[
  \left[
    C_{\mathcal{O}}^{(m)}(c),\;
    C_{\mathcal{O}}^{(m)}(c)+|\mathcal{U}|
  \right].
\]
Taking posterior quantiles of the lower and upper endpoints gives a 95\% posterior interval for the sample-analogue identified set. At the stated threshold, the interval is 87 to 231 LGAs, with width 144; 120 of that width reflects the assumption-free contribution of the unobserved LGAs. This interval is not an Imbens--Manski confidence set. It is a posterior interval for a sample-analogue identified set, conditional on the model for observed LGAs and on the population denominators used to calculate \(p\).

For the population share, let \(\pi_{\mathcal{O}}^{(m)}\) be the posterior share of people in observed LGAs who live in an area above the crisis threshold. Without assumptions about the unobserved areas, the corresponding national share lies in
\[
  \left[
    p\,\pi_{\mathcal{O}}^{(m)},\;
    p\,\pi_{\mathcal{O}}^{(m)}+(1-p)
  \right].
\]
The observed-only share is 13.72\%, and the identified interval is 10.63\% to 24.90\%. The lower endpoint is not a direct-data bound: it uses model-smoothed prevalences for observed LGAs. The upper endpoint allows every unobserved LGA to be above the threshold. We impose no monotonicity or other outcome restriction on \(\mathcal{U}\). The interval therefore displays both the model dependence at the observed end and the uncertainty that remains because no cluster was sampled in the unobserved areas.

\subsection{Spatial calibration protocol}

We evaluated area-level prediction with four spatial holdout designs. All four designs, their scoring rule, and the pass criteria were fixed before the runs; none of the criteria was revised after seeing the results. They address different prediction settings, from predicting one area among observed neighbours to predicting an entire region with its neighbouring areas removed. The exposure measure makes their results comparable: for each held-out LGA, count the neighbouring LGAs unavailable to the fit that predicts it, then average that count over held-out LGAs. Exposure is therefore the mean number of a held-out LGA's neighbours that no fit ever sees.

The whole-region design withheld all households and LGAs in one geopolitical region, refit the model without them, and predicted the withheld region. This design tests regional extrapolation because the held-out LGAs lose their local observations and much of their neighbouring information. The single-area design withheld an LGA while retaining its neighbours in the training data. The three area-level designs used 200 held-out LGAs. For the matched and exact-exposure designs, the same held-out areas were arranged into different prediction groups so that the mean exposure could be compared at different levels. The exact-exposure design targeted the exposure of the 120 LGAs without sampled clusters. The withheld-area lists and group assignments should be retained with the analysis scripts so that another analyst can reproduce the same folds.

For held-out LGA \(s\), let \(\widehat{\theta}^{\,M}_s\) and \(V^{M}_s\) be the model prediction and its posterior predictive variance. Let \(\widehat{\theta}^{\,D}_s\) be the direct unweighted proportion among held-out households, with design variance \(V^{D}_s\) estimated by cluster-robust linearisation using the sampled cluster as the primary sampling unit and the survey stratum in the variance calculation. The standardized error is
\[
  z_s =
  \frac{\widehat{\theta}^{\,M}_s-\widehat{\theta}^{\,D}_s}
       {\sqrt{V^{M}_s+V^{D}_s}}.
\]
Using the direct estimate's design variance avoids treating the held-out proportion as error-free. For each design we report the mean and standard deviation of \(z_s\), and coverage, defined as the share of held-out areas with \(\lvert z_s\rvert \le 1.96\). Coverage is deliberately \emph{not} the share of direct estimates falling inside the model interval: the direct estimate carries its own sampling error, so that weaker check under-covers by construction and would understate the calibration. The direct estimate and the model prediction target the same unweighted quantity.

The four designs and their outcomes are summarised in Table~\ref{tab:calibration}. The whole-region design removed 774 LGAs in one pass. The other three designs held out 200 LGAs each. The single-area design kept neighbours and had exposure 0.88. The matched over-exposed design had exposure 1.67, while the exact-exposure design had exposure 0.95. The whole-region design had exposure 5.4.

\begin{table}[htbp]
\centering
\begin{tabular}{p{0.30\textwidth}rrrr}
\hline
Design & Exposure & Mean \(z\) & SD \(z\) & Coverage \\
\hline
Whole region removed & 5.4 & \(-0.654\) & 1.648 & 0.775 \\
Single areas removed, neighbours kept & 0.88 & \(-0.052\) & 1.111 & 0.940 \\
Matched over-exposed areas & 1.67 & \(-0.068\) & 1.181 & 0.900 \\
Exact exposure of unsampled LGAs & 0.95 & \(-0.058\) & 1.178 & 0.890 \\
\hline
\end{tabular}
\caption{Spatial validation results. Exposure is the mean number of a held-out LGA's neighbours not seen by the fit used to predict it.}
\label{tab:calibration}
\end{table}

The pass rule required the mean \(z\) to lie between \(-0.15\) and \(0.15\), its standard deviation to lie between 0.85 and 1.15, and coverage to lie between 0.90 and 0.98. Under this fixed rule, the single-area design met all three criteria. The whole-region design did not meet any of them. The matched over-exposed design met the mean and coverage criteria but its standard deviation was above the pass band. At the exposure corresponding to the unsampled LGAs, coverage was 0.890 and the standard deviation of \(z\) was 1.178, so the exact-exposure design did not meet the coverage or spread criterion. Its group-level coverage ranged from 0.878 to 0.900.

The exact-exposure verdict is sensitive to one prediction group: omitting that group changes the pass/fail classification. Each prediction group contained 40 LGAs. The group-level range and this sensitivity preclude treating the pooled coverage as a precise estimate of a stable rate. The three designs with exposure between 0.88 and 1.67 do not show a monotone exposure pattern, so we do not fit or describe a coverage curve. The contrast supported by these designs is between single-area interpolation among sampled neighbours and whole-region extrapolation.

For intervals at unsampled LGAs, we used an inflation factor \(k=1.178\) and widened the model intervals by that factor. Coverage after widening was 0.940. We also assessed the factor out of fold, estimating it from all but one prediction group and checking coverage in the omitted group in turn; the resulting coverages were 0.875, 1.000, 0.925, 0.975, and 0.950. This check does not make the held-out covariates uncertain: covariate values were supplied as known in every fold. Nor does area-level validation establish calibration for the joint count \(C(c)\); the count interval depends on cross-area dependence and was not validated by these four designs.

\subsection{Computation and reproducibility}

We fitted the model in R 4.5.2 with INLA 25.10.19. The posterior functionals were calculated from 20,000 joint posterior draws using seed 2026-11-22. The same draws were used for the LGA exceedance probabilities, count functionals, and partial-identification endpoints; no additional likelihood evaluation was required for those summaries.

The repository scripts record the stages of the analysis. Cluster-to-LGA assignment is in \texttt{scripts/lga\_assign.R}; the main model fit is in \texttt{scripts/lga\_model.R}; the spatial calibration code includes \texttt{scripts/calibration\_spatial.R}; the identification calculations are in \texttt{scripts/identif\_bounds.R}; and the count identities are checked in \texttt{scripts/theory\_counts.R}. the matched- and exact-exposure folds are in \texttt{scripts/calibration\_matched\_exposure.R} and \texttt{scripts/calibration\_exact\_exposure.R}; the covariate surface is built by \texttt{scripts/lga\_covariates.R} and the covariate refit is \texttt{scripts/calibration\_covariates.R}. Every numerical result reported in the paper should be linked to the script and saved output that produced it. No analysis or model fit was run while drafting this section.
}{}
\IfFileExists{results.tex}{% Results section. Numerical claims are restricted to the verified facts in T24.
\section{Results}

\subsection{Descriptive picture}

The national weighted prevalence of severe food insecurity was 29.10\%, while the unweighted prevalence in the analysis sample was 27.74\%. These summaries refer to different estimands. The fitted likelihood uses the unweighted household records, so the unweighted proportion is the directly comparable sample summary; the weighted estimate remains the survey-weighted national result. We report both rather than treating the fitted surface as a replacement for the weighted prevalence. This distinction also applies to the area results below: the model describes probabilities for the unweighted sample quantity, while the weighted figure describes the national survey estimate.

Direct state estimates ranged from 4.6\% in Niger to 85.0\% in Yobe, with a median of 23.5\%. The range shows that the state summaries vary widely, and the median locates the middle of that distribution well below its upper endpoint. These state values provide a national-scale description, but they do not show how estimates vary among local government areas within a state. The LGA-level posterior summaries therefore add a finer geographic description rather than replacing the state estimates.

Affirmation was less common for more severe food experiences. The item-level affirmation rate fell monotonically from 71\% for worrying about food to 33.5\% for not eating for a whole day, with no inversions in the ordering. The monotone gradient is visible across the reported endpoints and was not interrupted by an item with a higher affirmation rate at a more severe point in the sequence. It describes the distribution of reported experiences; it is not an alternative prevalence definition for the area model.

The two spatial-dependence statistics differed by outcome measure. For severe food insecurity, Moran's \(I\) was 0.204, with permutation \(p=0.018\) from 9,999 draws. For the broader moderate-or-severe measure, Moran's \(I\) was \(-0.012\), with \(p=0.427\). Thus, the permutation result indicates spatial dependence for the severe measure, whereas the corresponding test did not detect spatial dependence for the broader measure. The second result is reported alongside the first because the measures do not yield the same spatial pattern.

The non-significant result should not be read as proof that the broader measure has no spatial association. It records that this test did not distinguish its statistic from the null distribution. The significant result for severe food insecurity is also a global summary: it does not imply that all locations share the same level of dependence. The state range and the LGA estimates describe geographic variation, while the two Moran statistics summarize different outcomes at a broader scale.

\subsection{Area estimates}

The posterior means for the 774 LGAs ranged from 0.017 to 0.918. Figure~\ref{fig:estimates} displays those means in sorted order with their intervals, so readers can see both the spread of the fitted estimates and their area-specific uncertainty. The displayed range is a range of model estimates, not a range of direct sample proportions. The intervals are needed to read the ordered means: two nearby point estimates need not have the same posterior uncertainty, and the plot preserves that distinction rather than reducing the distribution to its endpoints.

The endpoints alone do not describe the full distribution of area estimates. They do not show how many LGAs lie near the crisis threshold, nor do they convey how certain the model is about an individual area. The sorted display keeps those two features visible together. It should therefore be read as a collection of posterior area summaries, rather than as a ranking of observed household proportions.

Sample coverage was incomplete at LGA level. A sampled cluster was present in 654 LGAs, while 120 had no sampled cluster; those 120 areas contain 11.04\% of the population. For the unclustered LGAs, the model's posterior means ranged from 0.046 to 0.584, and their mean exceedance probability was 0.132. These are posterior predictions rather than direct observations. No sampled cluster contributes a local household outcome in those areas, so the reported means and exceedance probabilities are model output and should not be described as measured LGA prevalence.

The mean exceedance probability of 0.132 is a summary of posterior probabilities across the unclustered LGAs. It is not the share of households in those areas with severe food insecurity, and it is not an observed prevalence. The posterior mean and the exceedance probability answer different questions: one summarizes the predicted area prevalence, and the other summarizes the model's probability that an area crosses the chosen line. Both quantities for these areas remain predictions because the survey supplies no sampled cluster there.

The model summaries also distinguish the area and cluster variation represented in the fit. The reported LGA precision was 0.904 and the cluster precision was 4.784. The spatial mixing parameter \(\phi\) was 0.971, with a 95\% interval from 0.903 to 0.997. These parameter summaries accompany the area estimates but do not remove the difference between an LGA with observed households and an LGA whose estimate is produced without a sampled cluster.

The parameter summaries and the calibration results describe separate aspects of the fit. The mixing parameter and precisions summarize fitted variation, whereas the holdout results assess prediction against areas excluded from the corresponding fit. A parameter estimate within its reported interval is therefore not, by itself, a validation result for the area-specific predictions.

The model fit took 40--55 seconds, and its WAIC was 37,942.17 across runs. The posterior summaries reported here use 20,000 draws with seed 2026-11-22. These computational summaries identify the fit and draw set behind the reported quantities; they are not evidence that every area-specific interval has the same predictive performance. The calibration results below assess prediction under distinct holdout designs.

\begin{figure}[htbp]
  \centering
  \fbox{\parbox{0.88\linewidth}{\centering Figure output pending.}}
  \caption{Posterior mean prevalence for each of the 774 LGAs, ordered from lowest to highest, with the corresponding interval for each area. Output file: \todo{figure output filename}.}
  \label{fig:estimates}
\end{figure}

\subsection{The count}

At the threshold \(c=0.50\), the posterior expected number of LGAs above the threshold was \(E[C]=114.975\) of 774. The 95\% posterior interval was 101 to 130, with posterior standard deviation 7.383. This interval summarizes the model's joint posterior draws and is conditional on the fitted dependence structure; the marginal exceedance probabilities alone do not determine it. The feasible marginal-only range for the count's standard deviation is reported in the next subsection, where the distinction between the model interval and the information in the marginals is made explicit.

The variance decomposition separates the contribution of each LGA's exceedance probability from the contribution of pairwise joint exceedance. At \(c=0.50\), the independent term was 40.258 and the pairwise term was 14.248, giving a pairwise-term ratio of 0.354. The identity residual was \(-2.7\times 10^{-3}\). The exchangeable-areas comparator had a ratio of 0.187. These summaries show that the pairwise component is positive in the fitted model, while the comparator gives a smaller ratio under its alternative dependence structure.

The pairwise contribution is the part that is absent when the area indicators are treated as independent. Its positive value means that the fitted joint draws contain additional count variation beyond the independent term. The exchangeable-area ratio is a comparator for this same count summary; it should not be read as a decomposition into a separate causal contribution from geography. The reported ratios describe the two fitted dependence structures, not a test of whether one structure is correct.

The pairwise term has visible Monte Carlo variation. The split-half estimates were 13.37 and 15.12, a difference of approximately \(\pm 1.7\) around the reported value. We therefore report the interval and variance decomposition together, rather than treating the pairwise contribution as an exact constant. The split-half comparison concerns numerical variation in this posterior summary; it does not change the separate fact that the joint dependence structure is an assumption of the fitted model.

Posterior classification also leaves a group of areas close to the crisis line. Thirty-two LGAs had exceedance probability above 0.95, 509 were classified below the threshold, and 40 had probabilities between 0.40 and 0.60. The expected count changed with the area-level threshold: \(E[C]\) was 196.28 at \(c=0.40\) and 61.28 at \(c=0.60\), compared with 114.975 at \(c=0.50\). These three expected counts summarize threshold sensitivity; the variance decomposition in Table~\ref{tab:count} is reported at the central threshold.

\begin{table}[htbp]
  \centering
  \caption{Posterior count summaries across thresholds and variance decomposition at \(c=0.50\). The 95\% interval at \(c=0.50\) is conditional on the fitted joint dependence model and is not determined by the marginal exceedance probabilities alone.}
  \label{tab:count}
  \begin{tabular}{lrr}
    \toprule
    Threshold \(c\) & \(E[C]\) & 95\% posterior interval \\
    \midrule
    0.40 & 196.28 & 180--213 \\
    0.50 & 114.975 & 101--130 \\
    0.60 & 61.28 & 51--73 \\
    \midrule
    \multicolumn{3}{l}{Variance decomposition and comparator at \(c=0.50\)} \\
    Posterior standard deviation & \multicolumn{2}{r}{7.383} \\
    Independent variance term & \multicolumn{2}{r}{40.258} \\
    Pairwise variance term & \multicolumn{2}{r}{14.248} \\
    Pairwise-term ratio & \multicolumn{2}{r}{0.354} \\
    Exchangeable-areas comparator ratio & \multicolumn{2}{r}{0.187} \\
    Identity residual & \multicolumn{2}{r}{\(-2.7\times 10^{-3}\)} \\
    \bottomrule
  \end{tabular}
\end{table}

\subsection{What the marginals can and cannot pin down}

The marginal exceedance probabilities do not identify the variance of the count. By Proposition~\ref{prop:sharp}, for the 774 supplied marginals, the feasible range of \(\operatorname{Var}(C)\) is 0.0244 to 5492.63, corresponding to a standard-deviation range of 0.156 to 74.11. The lower and upper endpoints are attained by the cyclic and comonotone couplings, respectively. The fitted model's variance of 54.509 lies 0.99\% of the way up this feasible range. Thus, the \(E[C]\) value is fixed by the marginals, but the uncertainty around that count depends on how the exceedances are joined across areas.

The reported 95\% count interval of 101 to 130 is therefore a property of the fitted model's joint dependence structure, not a property of the data's marginal probabilities alone. The same marginals permit a count standard deviation anywhere from 0.156 to 74.11, so they cannot by themselves select the reported interval. The distinction is also visible in the tail probability: under the comonotone coupling, \(\Pr(C\geq 101)=0.476\); under the cyclic coupling it is 1.000; under the fitted model it is 0.984. These values use the same marginals but different joint couplings. The model interval must consequently be read as conditional on the dependence structure estimated by the model, not as an interval that the marginal data alone establish.

\begin{figure}[htbp]
  \centering
  \fbox{\parbox{0.88\linewidth}{\centering Figure output pending.}}
  \caption{The feasible range of \(\operatorname{Var}(C)\) given the area-specific marginal exceedance probabilities, with the fitted model's variance marked. Output file: \todo{figure output filename}.}
  \label{fig:feasible}
\end{figure}

\subsection{Partial identification}

Allowing the unclustered areas to contribute either zero or one to the threshold count gives an identified interval from 87 to 231 LGAs, with width 144. The observed-only count was 99.13, with an interval from 87 to 111. The full identified interval contains the model-based 101 to 130 interval, but the latter is conditional on the fitted joint dependence structure and is not determined by the marginals alone. Of the identified interval's width, 120 counts arise from the unrestricted contribution of the unclustered areas. The lower endpoint still uses model-smoothed prevalences for LGAs with sampled clusters, so the interval combines uncertainty about the unsampled areas with the fitted estimates for the observed areas.

The population-share bounds tell a different story from the count of areas because the unclustered LGAs do not contain an equal share of the population. The observed-only share of people living in a crisis LGA was 13.72\%, and the identified interval for the national share was 10.63\% to 24.90\%. The width of the population-share bound was 0.11037, equal to \(1-p\). This width describes the share of people whose crisis-area status is not fixed by the observed areas; it is not the fraction of geographic units without a cluster.

The count interval and the population-share interval should not be substituted for one another. The count gives every LGA one unit of weight, while the population share gives each area weight according to the people represented there. The unclustered areas account for a smaller share of people than their share of geographic units, so the population-share bounds are narrower than an area count alone would suggest. The reported endpoints keep the uncertainty attached to people distinct from the uncertainty attached to the number of places.

Under the monotonicity variant, the population-share interval was 0.1372 to 0.2324, a reduction in width of 1.3 percentage points. This variant narrows the interval relative to the unrestricted bound, but it relies on an additional ordering assumption. We report it as a sensitivity result rather than replacing the unrestricted interval. The count and population-share results are posterior intervals for sample-analogue identified sets, not frequentist confidence sets.

\begin{figure}[htbp]
  \centering
  \fbox{\parbox{0.88\linewidth}{\centering Figure output pending.}}
  \caption{The identified interval for the count compared with the model-based count interval. The model interval is conditional on the fitted joint dependence structure and is not identified by the marginal probabilities alone. Output file: \todo{figure output filename}.}
  \label{fig:identified}
\end{figure}

\subsection{Calibration}

The four calibration designs gave different results depending on whether prediction concerned a single area among observed neighbours or a whole region without its local observations. Table~\ref{tab:calibration} reports exposure, the mean and standard deviation of \(z\), coverage, and the result against each pre-declared criterion. The pass rule required mean \(z\) between \(-0.15\) and 0.15, standard deviation between 0.85 and 1.15, and coverage between 0.90 and 0.98. Each area-level design held out 200 LGAs, except the whole-region design; the area-level designs used groups of 40 for prediction.

\begin{table}[htbp]
  \centering
  \caption{Spatial calibration results. Exposure is the mean number of a held-out LGA's neighbours that were not seen by the fit used to predict it. The pre-declared pass bands were mean \(z\) in \([-0.15,0.15]\), standard deviation of \(z\) in \([0.85,1.15]\), and coverage in \([0.90,0.98]\).}
  \label{tab:calibration}
  \begin{tabular}{lrrrrl}
    \toprule
    Design & Exposure & Mean \(z\) & SD \(z\) & Coverage & Verdicts (mean/SD/coverage) \\
    \midrule
    Whole region removed & 5.4 & \(-0.654\) & 1.648 & 0.775 & Fail / Fail / Fail \\
    Single areas, neighbours kept & 0.88 & \(-0.052\) & 1.111 & 0.940 & Pass / Pass / Pass \\
    Matched over-exposed & 1.67 & \(-0.068\) & 1.181 & 0.900 & Pass / Fail / Pass \\
    Exact exposure of the 120 & 0.95 & \(-0.058\) & 1.178 & 0.890 & Pass / Fail / Fail \\
    \bottomrule
  \end{tabular}
\end{table}

The whole-region design failed all three pre-declared criteria. Its mean \(z\) of \(-0.654\) fell outside the allowed range, its standard deviation of 1.648 exceeded the upper limit, and its coverage of 0.775 fell below the lower limit. The failure is therefore not confined to coverage: the mean error and the spread of standardized errors also failed their respective checks. This design tests prediction of an area in a region that was wholly withheld, and the results do not support treating the resulting regional predictions as calibrated.

The single-area design met all three criteria, with mean \(z=-0.052\), standard deviation 1.111, and coverage 0.940. The matched over-exposed design passed the mean and coverage checks but failed the standard-deviation check at 1.181. The exact-exposure design, which matches the exposure of the 120 unclustered LGAs, passed the mean check but failed both the standard-deviation and coverage checks: its standard deviation was 1.178 and its coverage was 0.890. The exact-exposure result is therefore not a full calibration pass, even though its mean standardized error was within the pre-declared range.

The three area-level designs with exposure between 0.88 and 1.67 had coverages of 0.940, 0.890, and 0.900. These differences are within the noise of the held-out groups and do not support a monotone relationship between exposure and coverage. The supported contrast is between predicting a single area while its neighbours remain in the fit and predicting a whole region after its neighbours are removed. The whole-region coverage was 0.775, below the results for the three area-level designs, but the middle values should not be presented as a coverage curve.

At the exposure corresponding to the unclustered areas, group-level coverage ranged from 0.878 to 0.900. An interval inflation factor of \(k=1.178\) raised coverage to 0.940 after widening. The out-of-fold coverage values for the factor were 0.875, 1.000, 0.925, 0.975, and 0.950. These results report the behavior of intervals for areas under the stated prediction designs; they do not establish calibration for the joint count distribution.

The pooled coverage after widening is within the pre-declared coverage band, but the out-of-fold values vary across prediction groups. The exact-exposure design before widening still failed its coverage criterion, so the unadjusted result remains a failure. The widening result is a separate summary for area intervals and does not turn the original exact-exposure check into a pass. Neither result supplies a calibration test for the count formed from jointly dependent area indicators.

The correlation with direct estimates was 0.974 in sample, 0.768 out of sample under the single-area design, and 0.195 when a whole region was withheld. The in-sample and single-area results describe agreement when neighbouring areas remain represented in the fit, while the whole-region result describes a more demanding regional prediction. The gap between the latter two is consistent with the contrast in coverage, but neither area-level correlation validates the model-based count interval. That count interval remains conditional on the model's joint dependence structure, which the marginal probabilities alone do not identify.

\begin{figure}[htbp]
  \centering
  \fbox{\parbox{0.88\linewidth}{\centering Figure output pending.}}
  \caption{Observed coverage plotted against nominal coverage for the four spatial holdout designs. Output file: \todo{figure output filename}.}
  \label{fig:calibration}
\end{figure}

\subsection{Area covariates}

Three of the eight candidate area covariates could not be interpolated adequately. Travel time had an interpolation correlation of 0.614, livestock 0.535, and nightlights 0.522; all three were dropped as un-interpolatable. Five were retained. Population density was not interpolated at all: it was computed from population projections and polygon area. The other four were aridity at 0.999, elevation at 0.977, malaria at 0.923 and vegetation at 0.920. Precipitation interpolated just as well, at 0.999, but was dropped because it correlates with aridity at 0.985 and the selection rule keeps one of any such pair. Table~\ref{tab:covariates} reports all candidate correlations and their disposition. The distinction is between variables that could be carried to the area level and variables that could not, not between covariates with and without an association with the outcome.

\begin{table}[htbp]
  \centering
  \caption{Cross-validated interpolation correlations and disposition of the eight candidate area covariates.}
  \label{tab:covariates}
  \begin{tabular}{lrl}
    \toprule
    Candidate & Interpolation correlation & Disposition \\
    \midrule
    Population density & not interpolated & Kept \\
    Aridity & 0.999 & Kept \\
    Precipitation & 0.999 & Dropped: collinear with aridity, $r=0.985$ \\
    Elevation & 0.977 & Kept \\
    Malaria & 0.923 & Kept \\
    Vegetation & 0.920 & Kept \\
    Travel time & 0.614 & Dropped: un-interpolatable \\
    Livestock & 0.535 & Dropped: un-interpolatable \\
    Nightlights & 0.522 & Dropped: un-interpolatable \\
    \bottomrule
  \end{tabular}
\end{table}

The five interpolatable covariates were largely inert in the area model. After they were added, 19 of 20 coefficient intervals contained zero. Regional-design coverage increased from 0.775 to 0.797, but remained below the pre-declared lower limit of 0.90. The covariate refit therefore increased regional coverage without resolving the regional calibration failure, and it did not yield clear nonzero coefficient estimates for most included terms.
}{}
\IfFileExists{discussion.tex}{\section{Discussion}

The decision-relevant result is a range for the number of local government areas (LGAs) above the 50\% threshold, not a single point estimate. Severe food insecurity showed positive state-level spatial association (Moran's \(I=0.204\), permutation \(p=0.018\)), whereas the moderate-or-severe measure did not (\(I=-0.012\), \(p=0.427\)). This contrast is consistent with geographic structure concentrated in the severe tail rather than shared by the broader outcome. For agency use, we therefore report the identified range of 87 to 231 LGAs, not 115 as a fixed tally. The lower endpoint nevertheless depends on model-smoothed prevalences for observed LGAs, so it is not a direct-data count.

Under the fitted joint model, the posterior expected count was \(E[C]=114.975\) of 774 LGAs, with a 95\% posterior interval from 101 to 130 and a posterior standard deviation of 7.383. That interval describes a joint distribution of threshold crossings. The marginal exceedance probabilities alone permit \(\operatorname{Var}(C)\) from 0.0244 to 5492.63, equivalent to a standard-deviation range from 0.156 to 74.11; the fitted variance of 54.509 lies only 0.99\% of the way up this feasible range. Thus, the interval from 101 to 130 is a property of the fitted dependence structure, not something identified by the marginals. The marginals fix the expected count, but they do not say whether areas cross the line together. Area-specific decisions were also uneven: 32 LGAs had an exceedance probability above 0.95, 509 were classified below the line, and 40 lay in the 0.40--0.60 probability band.

The fitted dependence contribution is not negligible. The independent variance term was 40.258, while the pairwise term was 14.248, giving a pairwise-term ratio of 0.354. Under an exchangeable area effect, the ratio was 0.187. This comparison is consistent with roughly half of the pairwise contribution reflecting shared coefficient uncertainty and half reflecting geographic dependence. The pairwise estimate itself has Monte Carlo variation: split-half estimates were 13.37 and 15.12. The decomposition therefore helps explain why the count varies more than it would under independence, but it does not establish that the fitted spatial dependence is the only plausible joint structure.

The partial-identification result shows why the expected count should not be treated as an observed total. The survey includes at least one cluster in 654 LGAs, while 120 LGAs have none. Those unclustered areas are 15.5\% of all LGAs but contain 11.04\% of the population (\(p=0.88963\)). The model-based count over the observed 654 areas spans 87 to 111, while allowing each unclustered area to contribute either zero or one adds 120 LGAs to the width of the full 144-LGA interval. In other words, 120 of the 144-LGA width rests on an assumption about places without sampled households. The observed-only share of people living in a crisis LGA was 13.72\%, and the corresponding identified interval for the population share was 10.63\% to 24.90\%. These are different decision quantities: the area count gives every LGA equal weight, while the population share weights areas by their residents. The interval is a posterior interval for the sample-analogue identified set, conditional on the model for observed LGAs and the population denominators. Its lower endpoint is model-based rather than a direct-data bound.

The validation results determine how far the area estimates can be carried. When one LGA was withheld while its neighbours remained in the fit, the design met all three pre-declared criteria at exposure 0.82, with coverage of 0.940. The matched over-exposed design at exposure 1.67 had coverage of 0.900 but failed the spread criterion, with a standardized-error standard deviation of 1.181. The whole-region holdout failed all three criteria at exposure 5.4, with coverage of 0.775. Agreement with direct estimates also depended on the holdout: the correlation was 0.975 in sample, 0.768 under the single-area design, and 0.195 when a whole region was withheld. In the latter design, predictions remained near 0.13 to 0.28 across areas with different observed outcomes. The contrast supports interpolation among sampled neighbours, not extrapolation to a region without local observations. High in-sample agreement therefore does not justify treating the fitted surface as portable across regions.

For the 120 unclustered LGAs, the exact-exposure design is the more relevant check. At exposure 0.95, unadjusted coverage was 0.890 and the standard deviation of standardized errors was 1.178, so the design failed its coverage and spread criteria. Group-level coverage ranged from 0.878 to 0.900, and one group of 40 LGAs carried the pass-or-fail verdict. Widening an individual-area interval by \(k=1.178\) raised pooled coverage to 0.940. The out-of-fold coverages for that widening factor were 0.875, 1.000, 0.925, 0.975, and 0.950. An agency using an interval for an unclustered LGA should report the widened interval and describe its observed coverage in these tests, about 0.89 before widening and 0.94 after. It should not label that interval a nominal 95\% interval. The variation across groups also means that the pooled coverage should not be read as a stable guarantee for every unsampled area.

The area-covariate analysis is informative because it separates a lack of measurable signal from a lack of usable geographic data. Three of eight candidates could not be interpolated adequately from the 1,380 sample points: nightlights had a cross-validated interpolation correlation of 0.522, travel time 0.614, and livestock 0.535. The five covariates that could be carried to the area level were largely inert after household composition was accounted for: 19 of 20 coefficient intervals contained zero. These results do not show that economic activity, market access, or pastoral livelihoods are unrelated to food insecurity. Rather, the available measurements of those dimensions could not be recovered at this resolution, while the area variables that were recoverable added no clear signal in this model. A future analysis would need better geographic measurements before interpreting a weak area-level association as evidence against these mechanisms.

The literature comparison is narrow by design. Chacón-Montalván et al. (2025) report threshold-exceedance probabilities and hotspot counts, so neither quantity is presented here as a new use of spatial food-insecurity models. Their paper maps a continuous surface in one municipality, with eight census sectors as its finest administrative unit. Nigerian estimates already exist for single local government areas from cross-sectional surveys (Omuemu, Otasowie and Onyiriuka, 2012; Kayode and Alabi, 2025) and for the 109 senatorial districts from survey indices with hotspot classification (Jerumeh, 2024); we found in that work no interval for an area estimate and no interval for the count. What we did not find, and what this paper supplies, is a set of estimates for the national set of 774 local government areas together with an interval for the number of areas above a threshold. Related estimators for area-level proportions include bivariate spatial models for small-area proportions (Yu, Xu and Yu, 2025) and direct-assisted unit-level models for rare-event prevalence (McGovern, Wilson and Wakefield, 2026); in what we could read of this work we found no interval for the number of areas above a line. The contribution here is the sharp bound for the count, the pre-declared spatial validation protocol, and the identification statement for areas without a sampled cluster, applied to Nigerian LGAs; it is not a new estimator.

Several limits define what the results do not claim. The model provides no validated interval for the count of areas above the threshold. The whole-region holdout does not support regional extrapolation. The 120 LGAs without clusters do not have a calibrated map, and the likelihood does not use survey design weights. The area-level calibration tests do not repair these gaps: they evaluate predictions for particular holdout designs, not the joint count distribution or the unobserved prevalence in every unsampled LGA. The 120 areas cannot be scored against a known truth. The lower endpoint of the identification interval also depends on model-smoothed prevalence in observed areas, so the interval does not remove model dependence from the observed part of the count.

Two further sources of uncertainty affect the population interpretation. LGA assignment used geoBoundaries polygons rather than the survey's own administrative boundaries, and population denominators are 2020 projections based on the 2006 census. The main inferential limitation, however, is the unweighted likelihood. The unweighted national prevalence was 27.74\%, compared with the survey-weighted estimate of 29.10\%, a difference of about 1.4 percentage points. The LGA surface and its count therefore describe the unweighted sample quantity, not a design-weighted population prevalence.
}{}
\IfFileExists{conclusions.tex}{\section{Conclusions}

A count of local government areas (LGAs) above a prevalence threshold is a distinct posterior quantity, not a list of area estimates. Its variance depends on pairwise exceedance probabilities, which marginal-only reporting does not supply. At the 50\% threshold, the fitted expected count was 115 of 774 LGAs, with a 95\% posterior interval of 101--130. The pairwise term means that an independence calculation would understate the fitted count's standard deviation by 14\%. This interval depends on the fitted joint dependence structure and was not validated by the area-level cross-validation.

The marginal probabilities allow the count's standard deviation to range from 0.156 to 74.11. The model's variance sits only 0.99\% of the way through the feasible variance range. The posterior interval for the sample-analogue identified count runs from 87 to 231. Of its 144-area width, 120 reflects unrestricted contributions from LGAs without sampled clusters. Those LGAs contain 11.04\% of the population, and the lower endpoint for sampled areas still uses model-smoothed prevalence estimates.

Spatial cross-validation used pass rules fixed before fitting. At the exposure faced by unclustered LGAs, area-level coverage was 0.890 against the nominal 0.95; withholding a whole region yielded 0.775. An interval-widening factor \(k=1.178\) was checked out of fold for area-level intervals. The protocol is reusable because it replaces an assumed nominal level with measured performance in prediction settings that match intended use.

The transferable parts are the sharp bounds and the pre-declared protocol, not a new estimator. The Nigerian map is a demonstration, and its area-level precision is supported only for prediction settings represented by sampled data.
}{}
\IfFileExists{declaration.tex}{\section{Declaration}

\noindent\textbf{Competing interests.} The author declares no competing interests.

\medskip
\noindent\textbf{Funding.} No funding was received for this work.

\medskip
\noindent\textbf{Ethics.} This study is a secondary analysis of the 2024 Nigeria Demographic and Health Survey. The DHS Program collects and releases these data under its own consent and anonymisation procedures. The de-identified files used here contain no identifiable information, and no further ethical approval was required for this analysis.

\medskip
\noindent\textbf{Data availability.} Household microdata and geographic files are available from the DHS Program to registered users and cannot be redistributed by the author. Code producing every result is released under the MIT licence. Derived results and documentation are released under CC BY 4.0 at \url{https://github.com/batestguy/food-insecurity-nigeria}. The analysis runs in R 4.5.2 with INLA 25.10.19.

\medskip
\noindent\textbf{Acknowledgement.} The author thanks the DHS Program for access to the data, and the National Bureau of Statistics and UNFPA for the population projections used as denominators.
}{}

% Theory section: docs/theory.tex is written as a set of propositions. It sits
% inside Methods unless the journal prefers a standalone section; that decision
% is recorded in PLAN.md.

\begin{sloppypar}
\raggedright
\begin{thebibliography}{99}
\bibitem{cafiero2024}
Cafiero, C., Gheri, F., \& Viviani, S. (2024). Validating the food insecurity experience scale for
use in analyses of recent food insecurity. \emph{Global Food Security, 42}, 100783.
\url{https://doi.org/10.1016/j.gfs.2024.100783}

\bibitem{chacon2025}
Chacón-Montalván, E. A., Parry, L., Giorgi, E., Torres, P., Orellana, J. D. Y., Moraga, P., \&
Taylor, B. M. (2025). Mapping food insecurity in the Brazilian Amazon using a spatial item factor
analysis model. \emph{The Annals of Applied Statistics, 19}(4), 3438--3463.
\url{https://doi.org/10.1214/25-aoas2072}

\bibitem{chandra2022}
Chandra, H., \& Verma, B. (2022). Small area estimation of food insecurity prevalence for the state
of Uttar Pradesh in India. \emph{Model Assisted Statistics and Applications, 17}(2), 73--85.
\url{https://doi.org/10.3233/mas-220011}

\bibitem{jerumeh2024}
Jerumeh, T. R. (2024). Mapping the spatial distribution of household food insecurity in Nigeria.
\emph{GeoJournal, 89}(4), Article 154. \url{https://doi.org/10.1007/s10708-024-11035-4}

\bibitem{kayode2025}
Kayode, O. O., \& Alabi, Q. K. (2025). Prevalence and contributing factors of household food
insecurity among women in Akinyele Local Government, Ibadan, Oyo State, Nigeria. \emph{Pan African
Medical Journal, 52}, 50. \url{https://doi.org/10.11604/pamj.2025.52.50.36958}

\bibitem{macbride2025}
MacBride, C., Davies, V., \& Lee, D. (2025). A spatial autoregressive random forest algorithm for
small-area spatial prediction. \emph{The Annals of Applied Statistics, 19}(1), 485--504.
\url{https://doi.org/10.1214/24-aoas1969}

\bibitem{mcgovern2026}
McGovern, A., Wilson, K., \& Wakefield, J. (2026). Direct-assisted Bayesian unit-level modeling for
small area estimation of rare event prevalence. \emph{Journal of Survey Statistics and Methodology,
14}(1), 178--208. \url{https://doi.org/10.1093/jssam/smaf037}

\bibitem{michal2024}
Michal, V., Schmidt, A. M., Freitas, L. P., \& Cruz, O. G. (2024). A Bayesian hierarchical model for
disease mapping that accounts for scaling and heavy-tailed latent effects. \emph{Statistical Methods
in Medical Research, 34}(2), 307--321. \url{https://doi.org/10.1177/09622802241293776}

\bibitem{moura2025}
Moura, F. A. S., Pereira, S., \& Silva, G. L. (2025). A Beta--Beta Prime model for rates and their
precision for small area estimation: An application to Brazilian food insecurity index.
\emph{Journal of Survey Statistics and Methodology, 14}(2), 430--451.
\url{https://doi.org/10.1093/jssam/smaf002}

\bibitem{omuemu2012}
Omuemu, V. O., Otasowie, E. M., \& Onyiriuka, U. (2012). Prevalence of food insecurity in Egor local
government area of Edo State, Nigeria. \emph{Annals of African Medicine, 11}(3), 139.
\url{https://doi.org/10.4103/1596-3519.96862}

\bibitem{osung2026}
Osung, P. O. (2026). Household-level food-insecurity prediction from experience-scale surveys and
open geospatial data: A reproducible machine-learning pipeline for Nigeria. \emph{PLOS One, 21}(10),
e0349978. \url{https://doi.org/10.1371/journal.pone.0349978}

\bibitem{riebler2016}
Riebler, A., Sørbye, S. H., Simpson, D., \& Rue, H. (2016). An intuitive Bayesian spatial model for
disease mapping that accounts for scaling. \emph{Statistical Methods in Medical Research, 25}(4),
1145--1165. \url{https://doi.org/10.1177/0962280216660421}

\bibitem{rue2009}
Rue, H., Martino, S., \& Chopin, N. (2009). Approximate Bayesian inference for latent Gaussian models
by using integrated nested Laplace approximations. \emph{Journal of the Royal Statistical Society:
Series B (Statistical Methodology), 71}(2), 319--392.
\url{https://doi.org/10.1111/j.1467-9868.2008.00700.x}

\bibitem{santillana2023}
Santillana, M. A., Gallegos, J. A., García, A. I., Díaz, E., Gutiérrez, D., \& González, N. L.
(2023). Perception of insecurity in municipalities in Mexico: A small area estimation approach.
\emph{Statistical Journal of the IAOS, 39}(2), 491--500. \url{https://doi.org/10.3233/sji-220101}

\bibitem{smalley2024}
Smalley, H., \& Edwards, K. (2024). Chronic back pain prevalence at small area level in England: The
design and validation of a 2-stage static spatial microsimulation model. \emph{Spatial and
Spatio-temporal Epidemiology, 48}, 100633. \url{https://doi.org/10.1016/j.sste.2023.100633}

\bibitem{villacis2022}
Villacis, A. H., Mayorga, J., \& Mishra, A. K. (2022). Experience-based food insecurity and
agricultural productivity in Nigeria. \emph{Food Policy, 113}, 102286.
\url{https://doi.org/10.1016/j.foodpol.2022.102286}

\bibitem{wulandari2025}
Wulandari, I., Darmawan, E. S., \& Zamzanah, S. (2025). Spatial Bayesian small area estimation of
stunting prevalence at the subdistrict level. \emph{InPrime: Indonesian Journal of Pure and Applied
Mathematics, 7}(2), 128--143. \url{https://doi.org/10.15408/inprime.v7i2.46491}

\bibitem{yilema2022}
Yilema, S. A., Shiferaw, Y. A., Zewotir, T. T., \& Muluneh, E. K. (2022). Spatial small area
estimates of undernutrition for under five children in Ethiopia via combining survey and census data.
\emph{Spatial and Spatio-temporal Epidemiology, 42}, 100509.
\url{https://doi.org/10.1016/j.sste.2022.100509}

\bibitem{yoo2024}
Yoo, M., \& Koo, H. (2024). Exploring the impact of spatial autocorrelation on optimistic bias in
cross-validation and assessing the effectiveness of spatial cross-validation. \emph{Cartography and
Geographic Information Science, 52}(5), 596--609. \url{https://doi.org/10.1080/15230406.2024.2422593}

\bibitem{yu2025}
Yu, H., Xu, X., \& Yu, L. (2025). Hierarchical Bayesian bivariate spatial modeling of small area
proportions with application to health survey data. \emph{Statistical Methods in Medical Research,
34}(5), 1018--1036. \url{https://doi.org/10.1177/09622802251316968}
\end{thebibliography}
\end{sloppypar}

\end{document}
